\documentclass[../Main.tex]{subfiles}

\begin{document}
\section{Scaling and Equations}
Consider a general thin layer flow, as in figure~\ref{figThinLayer}.
\begin{figure}
    \centering
    \begin{tikzpicture}[scale=1]
        \draw[hatch side] (0, 0) -- (3, 0);

    \end{tikzpicture}
    \caption{Diagram of Thin Layer Flow}
    \label{figThinLayer}
\end{figure}
In particular, we need a long and thin geometry: $h \ll L, \frac{\partial h}{\partial x} \ll 1$. To build the governing equations we need the following ingredients:

\begin{enumerate}
    \item Applying $\nabla \cdot \vec{u} = 0$, $\frac{\partial u}{\partial x} + \frac{\partial v}{\partial y} = 0$ which means $\frac u L \sim \frac v h$, and so $v \sim \frac h L u \ll u$. That is, fluid is almost horizontal -- quasi-parallel to the boundary.
    \item In $\mu \nabla^2\vec u$ we can approximate $\nabla^2$ by $\frac{\partial^{2}}{\partial y^{2}}$. This is because $h \ll L$ and so $\frac{\partial}{\partial y} \gg \frac{\partial}{\partial x}$.
\end{enumerate}
\end{document}